\begin{tcolorbox}[
    colback=promptbg!100!white,
    sharp corners=south, 
    boxrule=0.25mm, 
    fonttitle=\bfseries, 
    title={\textbf{Orchestrator: topology sampling}}, 
    width=\textwidth, 
    enhanced,
    drop shadow
    ]


You are an orchestrator managing a team of specialized agents for multi-hop 
question answering. Your goal is to design an effective multi-agent execution topology 
to answer the given query correctly. \\[2mm]
\textbf{You have access to the following agents:} \\
\{agent\_descriptions\}

\vspace{0.5em}
\textbf{You have retrieved the following relevant experiences from past executions:} \\
\{retrieved\_experiences\}

\vspace{0.5em}
\textbf{Each experience entry has the format:}
\begin{itemize}
    \item Query Type: \{query\_type\}
    \item Insight: \{insight\}
    \item Utility score: \{utility\}
\end{itemize}

\vspace{0.5em}
\textbf{Given the query below, design a coordination topology by:}
\begin{enumerate}
    \item Selecting the subset of agents best suited to this query type.
    \item Specifying their execution order (sequential or parallel where appropriate).
    \item Defining dependency relationships (which agent's output feeds into which).
\end{enumerate}

\vspace{0.5em}
\textbf{Query:} \{query\}

\vspace{0.5em}
\textbf{Respond in the following JSON format:}
\begin{verbatim}
{
  "query_profile": "<one-sentence characterization of query type>",
  "selected_agents": ["<agent_name>", ...],
  "execution_order": [
    {"step": 1, "agent": "<agent_name>", "depends_on": [],
    
    "mode": "sequential|parallel"},
    ...
  ]
}
\end{verbatim}
\end{tcolorbox}

\begin{tcolorbox}[colback=promptbg!100!white,title=Orchestrator: Semantic Advantage Extraction]
\small
You are an orchestrator evaluating a group of multi-agent execution trajectories
for the same query. Your goal is to identify why some trajectories succeeded and 
others failed, and extract generalizable insights to guide future orchestration.

\vspace{0.5em}
\textbf{Query:} \{query\} \\
\textbf{Query type:} \{query\_type\}

\vspace{0.5em}
The following \{G\} trajectories were executed, ranked by task performance (F1) 
and then by efficiency (total tokens consumed):

\vspace{0.5em}
\{trajectory\_group\}

\vspace{0.5em}
\textbf{Each trajectory entry includes:}
\begin{itemize}
    \item Topology used (agents selected, execution order)
    \item Per-agent actions, inputs and outputs
    \item Final answer produced
    \item F1 score: \{f1\}
    \item Total tokens consumed: \{tokens\}
\end{itemize}

\vspace{0.5em}
\textbf{Perform a comparative analysis across the group:}
\begin{enumerate}
    \item Identify what the successful trajectories did differently from the failed ones.
    \item Identify which agents, orderings, or dependencies contributed to success or failure.
    \item Identify recurring failure patterns (e.g., incorrect agent selection, redundant steps, missing retrieval before reasoning).
    \item Distill findings into concise, actionable insights applicable to future queries of the same type.
\end{enumerate}

\vspace{0.5em}
\textbf{Respond in the following JSON format:}
\begin{verbatim}
{
  "success_factors": ["<factor>", ...],
  "failure_modes": ["<failure_mode>", ...],
  "insights": [
    {
      "query_type": "<type this insight applies to>",
      "insight": "<actionable natural language insight>"
    },
    ...
  ]
}
\end{verbatim}
\end{tcolorbox}





\begin{tcolorbox}[
    colback=promptbg!100!white,
    sharp corners=south, 
    boxrule=0.25mm, 
    fonttitle=\bfseries, 
    title={\textbf{Experience Library Operations)}}, 
    width=\textwidth, 
    enhanced,
    drop shadow
    ]
You are managing an experience library for a multi-agent question answering system. 
The library stores insights as structured entries of the form:
\[
(\text{query\_profile } c,\ \text{insight } z,\ \text{utility } u)
\]
where utility $u$ reflects how often insight $z$ has led to successful agent orchestration.

\vspace{0.5em}
\textbf{Current experience library:} \\
\{current\_library\}

\vspace{0.5em}
\textbf{New insights extracted from the latest execution group:} \\
\{new\_insights\}

\vspace{0.5em}
\textbf{For each new insight, decide the appropriate consolidation operation:}
\begin{itemize}
    \item \textbf{ADD}: Insert a distinct insight not covered by existing entries.
    \item \textbf{MERGE}: Combine semantically similar or complementary insights into a single, more complete entry.
    \item \textbf{PRUNE}: Remove a lower-utility entry that conflicts with a higher-utility one.
    \item \textbf{KEEP}: Retain the current library without modification.
\end{itemize}

\vspace{0.5em}
\textbf{Guidelines:}
\begin{itemize}
    \item Prefer \textbf{MERGE} over \textbf{ADD} when insights share the same query type and recommend compatible strategies.
    \item Prefer \textbf{PRUNE} when entries provide contradictory guidance for the same query type and differ substantially in utility.
    \item Avoid unbounded growth—consolidate aggressively to maintain generalizability.
\end{itemize}

\vspace{0.5em}
\textbf{Respond in the following JSON format:}
\begin{verbatim}
{
  "operations": [
    {
      "operation": "ADD|MERGE|PRUNE|KEEP",
      "new_insight": "<text of new insight>",
      "target_entry_ids": ["<id of existing entry if MERGE or PRUNE>"],
      "merged_insight": "<combined insight text if MERGE, else null>",
      "rationale": "<one sentence explaining the decision>"
    },
    ...
  ]
}
\end{verbatim}
\end{tcolorbox}


\begin{tcolorbox}[
    colback=promptbg!100!white,
    sharp corners=south, 
    boxrule=0.25mm, 
    fonttitle=\bfseries, 
    title={\textbf{Prompt for Agent Prompt Integration}}, 
    width=\textwidth, 
    enhanced,
    drop shadow
    ]

\textbf{System Prompt:} \\
You are responsible for managing and evolving the prompt of a single agent under strict length and structural constraints. Your task is to refine, optimize, and consolidate the agent’s prompt to ensure clarity, consistency, and adherence to operational and behavioral requirements. \\[2mm]

\textbf{User Prompt:} \\

\textbf{Current Agent Prompt:} \\
\begin{verbatim}
{CURRENT_PROMPT}
\end{verbatim}

\textbf{Proposed Operational Rules:} \\
\begin{verbatim}
{operational_rules}
\end{verbatim}

\textbf{Proposed Behavioral Principles:} \\
\begin{verbatim}
{behavioral_principles}
\end{verbatim}

\textbf{Constraints:} 
\begin{enumerate}
    \item Limit the number of tactical rules to a maximum of \textbf{{K}}.  
    \item All instructions must be internally consistent — no contradictions.
    \item Preserve the agent's core role definition and tool usage instructions
    \item Remove redundant, contradictory, or ambiguous instructions.  
    \item Preserve essential operational and behavioral requirements.  
    \item Ensure the updated prompt is concise, coherent, and actionable.  
\end{enumerate}

\textbf{Task:} \\
Produce a \textbf{prompt diff} that clearly indicates the modifications required to integrate the proposed rules and principles into the current agent prompt while satisfying the constraints above. Highlight additions, deletions, and replacements in a structured format.

\end{tcolorbox}


\begin{tcolorbox}[
    colback=promptbg!100!white,
    sharp corners=south, 
    boxrule=0.25mm, 
    fonttitle=\bfseries, 
    title={\textbf{RoPE: contrastive analysis}}, 
    width=\textwidth, 
    enhanced,
    drop shadow
    ]

You are analyzing the results of prompt variant re-executions for a failing agent to extract concrete prompt improvements.

\textbf{Agent name:} \{agent\_name\} \\
\textbf{Agent role:} \{agent\_role\} \\
\textbf{Original prompt:} \{original\_prompt\}

\textbf{Variant execution results:} \\
\{variant\_results\}

Each result entry contains:
\begin{itemize}
  \item \textbf{Variant prompt:} \{variant\_prompt\}
  \item \textbf{Re-executed trajectory:} \{trajectory\}
  \item \textbf{F1 score:} \{f1\}
  \item \textbf{Tokens used:} \{tokens\}
\end{itemize}

Perform a contrastive analysis between successful and failed variants:

\begin{enumerate}
  \item \textbf{Operational rules} ($\Delta \rho_i^{op}$): Extract short-term corrective behaviors — specific, concrete instructions that directly address the observed failure pattern. These should be actionable in the agent's very next execution. \\

  \item \textbf{Behavioral principles} ($\Delta \rho_i^{bp}$): Extract long-term strategic generalizations — higher-level guidance distilled from patterns across multiple trajectories. These shape the agent's overall approach rather than fixing a single failure. \\

\end{enumerate}
Respond in the following JSON format:
\begin{verbatim}
{
  "operational_rules": [
    {
      "rule": "<concrete instruction>",
      "derived_from": "<which variant comparison motivated this rule>"
    },
    ...
  ],
  "behavioral_principles": [
    {
      "principle": "<strategic guidance>",
      "derived_from": "<pattern observed across which variants>"
    },
    ...
  ],
  "updated_prompt": "<full revised prompt integrating all rules and principles \\ 
  into the original, with redundant instructions pruned>"
}
\end{verbatim}
\end{tcolorbox}